%==============================================================================
% Structural Change Cheatsheet
% Lecture 6 — Macroeconomics I (BSE)
% Exam-Focused Reference
%==============================================================================

\documentclass[a4paper,11pt]{article}

\usepackage{geometry}
\geometry{margin=1in, headheight=15pt}

\usepackage{amsmath}
\usepackage{amssymb}
\usepackage{mathtools}
\usepackage{graphicx}

\usepackage{xcolor}
\definecolor{blue3}{RGB}{0,0,180}
\definecolor{green3}{RGB}{0,120,0}
\definecolor{orange3}{RGB}{200,100,0}

\usepackage{siunitx}
\sisetup{detect-all}

\usepackage[colorlinks=true,linkcolor=blue3,citecolor=blue3,urlcolor=blue3]{hyperref}

\usepackage{cleveref}

\usepackage{tcolorbox}
\tcbuselibrary{skins,breakable}

\newtcolorbox{derivationbox}[1][]{colback=blue!10,colframe=blue3,coltitle=blue3,title={Derivation:},fonttitle=\bfseries,#1}
\newtcolorbox{resultbox}[1][]{colback=green!10,colframe=green3,coltitle=green3,title={Result Summary:},fonttitle=\bfseries,#1}
\newtcolorbox{keyformulabox}[1][]{colback=orange!10,colframe=orange3,coltitle=orange3,title={Key Formula:},fonttitle=\bfseries,#1}

\usepackage{booktabs}
\usepackage{enumitem}
\usepackage{tikz}

% Sector notation commands
\newcommand{\LA}{L_A}
\newcommand{\LM}{L_M}
\newcommand{\LS}{L_S}
\newcommand{\YA}{Y_A}
\newcommand{\YM}{Y_M}
\newcommand{\YS}{Y_S}

\title{\textbf{Structural Change}\\[0.3em]
       \large Lecture 6 Cheatsheet --- Macroeconomics I (BSE)}
\author{Exam Reference}
\date{}

\begin{document}

\maketitle
\tableofcontents
\clearpage
\pagestyle{plain}

%==============================================================================
\section{Notation and Definitions}
%==============================================================================

\subsection*{Sector Variables}
\begin{align*}
& A = \text{Agriculture}, \quad M = \text{Manufacturing}, \quad S = \text{Services}\\
& c_t^i = \text{per-capita consumption in sector } i \in \{A, M, S\}\\
& L_{it} = \text{employment in sector } i & Y_{it} &= \text{output in sector } i\\
& p_i = \text{price of good } i & I &= \text{aggregate income/expenditure}\\
& \gamma^A, \, \gamma^S = \text{Stone-Geary subsistence parameters} (\gamma^A, \gamma^S > 0)\\
& \eta_A, \, \eta_M, \, \eta_S = \text{expenditure-share parameters}, \quad \eta_A + \eta_M + \eta_S = 1\\
& g = \text{aggregate consumption growth rate (BGP)} & n &= \text{population growth rate}\\
& B_i = \text{sector-}i \text{ productivity level} & A_{it} &= \text{TFP of sector } i\\
& a_i = \dot{A}_{it}/A_{it} = \text{TFP growth rate in sector } i\\
& \varepsilon = \text{elasticity of substitution between goods (CES aggregator)}\\
& \alpha_i = \text{capital share of sector } i & k_t &= K_{1t}/K_t = \text{capital allocation fraction}\\
& \lambda_t = L_{1t}/L_t = \text{labor allocation fraction (sector 1)}\\
& \sigma = \text{intertemporal elasticity of substitution (HRV preferences)}\\
& \bar{c}_i = \text{non-homothetic threshold parameter (HRV)}
\end{align*}

\subsection*{Key Notation Conventions}
\begin{itemize}[nosep]
    \item Dot notation: $\dot{x}_t \equiv dx_t/dt$, growth rate $\gamma_x \equiv \dot{x}/x$.
    \item Subscript $t$ on time-varying variables may be dropped when clear from context.
    \item ``Sector 1'' and ``Sector 2'' used in the supply model; agriculture/manufacturing/services in demand model.
    \item Engel curve: relationship between income and quantity demanded holding prices fixed.
\end{itemize}

\clearpage

%==============================================================================
\section{Stylized Facts of Structural Change}
%==============================================================================

\subsection*{Cross-Country Patterns (as Income Rises)}

\begin{center}
\begin{tabular}{lccc}
\toprule
\textbf{Sector} & \textbf{Employment Share} & \textbf{Output Share} & \textbf{Pattern} \\
\midrule
Agriculture & $\LA/L$ & $\YA/Y$ & Declines monotonically \\
Manufacturing & $\LM/L$ & $\YM/Y$ & Inverse-U (rises then falls) \\
Services & $\LS/L$ & $\YS/Y$ & Rises monotonically \\
\bottomrule
\end{tabular}
\end{center}

\subsection*{Productivity Growth Rates (Duarte \& Restuccia, 1956--2004)}

\begin{center}
\begin{tabular}{lc}
\toprule
\textbf{Sector} & \textbf{Average TFP Growth Rate} \\
\midrule
Agriculture & 3.8\% per year \\
Manufacturing & 2.4\% per year \\
Services & 1.3\% per year \\
\bottomrule
\end{tabular}
\end{center}

\begin{resultbox}
\textbf{Key observation:} TFP growth rates differ systematically across sectors. Agriculture and manufacturing have high productivity growth; services have low productivity growth. Relative price changes (services become relatively expensive) and relative income effects both drive reallocation of labor.
\end{resultbox}

\subsection*{Two Candidate Mechanisms}

\begin{enumerate}
    \item \textbf{Demand-side (non-homotheticities):} As income grows, income elasticities differ across sectors. People spend a declining share on food (necessity) and a rising share on services (luxury). Even at constant relative prices, sector shares change.
    \item \textbf{Supply-side (Baumol cost disease):} Productivity growth differs across sectors. Fast-growing sectors have falling relative prices, inducing substitution. If goods are complements ($\varepsilon < 1$), the slow-growing sector absorbs factors.
\end{enumerate}

\clearpage

%==============================================================================
\section{Demand-Side: Homotheticity and Engel Curves}
%==============================================================================

\subsection*{Homothetic Preferences}

Preferences are \textbf{homothetic} if the utility function is homogeneous of degree 1:
\[
U(\lambda c^A, \lambda c^M, \lambda c^S) = \lambda\, U(c^A, c^M, c^S) \quad \forall\, \lambda > 0
\]

\textbf{Implication for Engel curves:} With homothetic preferences, all Engel curves are linear \emph{and pass through the origin}. Doubling income exactly doubles demand for every good. Sectoral budget shares are independent of income.

\subsection*{Engel Curves}

The Engel curve for good $i$ is $x_i = x_i(I, p_i)$ holding prices fixed. Classification:

\begin{center}
\begin{tabular}{lll}
\toprule
\textbf{Type} & \textbf{Income Elasticity $\varepsilon_I$} & \textbf{Engel Curve Shape} \\
\midrule
Normal good & $\varepsilon_I > 0$ & Increasing in $I$ \\
Necessity & $0 < \varepsilon_I < 1$ & Increasing and concave in $I$ \\
Luxury & $\varepsilon_I > 1$ & Increasing and convex in $I$ \\
Inferior good & $\varepsilon_I < 0$ & Decreasing in $I$ \\
Homothetic & $\varepsilon_I = 1$ & Linear through origin \\
\bottomrule
\end{tabular}
\end{center}

\subsection*{Engel's Law}

As per capita income rises, the \emph{budget share} of food (agriculture) declines. This is the oldest stylized fact in economics (Ernst Engel, 1857). Formally: food has income elasticity less than 1, so $\partial (p_A c^A / I)/\partial I < 0$.

\begin{resultbox}
\textbf{Non-homotheticity is required for structural change via demand.} If preferences were homothetic, income growth would not change sectoral budget shares, and no demand-driven reallocation would occur. The key is that agriculture is a necessity ($\varepsilon_I^A < 1$) and services tend to be a luxury ($\varepsilon_I^S > 1$).
\end{resultbox}

\clearpage

%==============================================================================
\section{Stone-Geary Preferences and Demand Derivation}
%==============================================================================

\subsection*{Stone-Geary Utility}

The Stone-Geary utility function introduces subsistence levels:
\begin{align*}
u_t = (c_t^A - \gamma^A)^{\eta_A} \cdot (c_t^M)^{\eta_M} \cdot (c_t^S + \gamma^S)^{\eta_S}
\tag{Stone-Geary}
\end{align*}

\noindent where $\gamma^A > 0$ (subsistence floor for agriculture), $\gamma^S > 0$ (satiation level for services with reversed sign), and $\sum_{i} \eta_i = 1$, $\eta_i > 0$.

General compact form:
\[
u = \prod_{i \in \{A,M,S\}} (c_i - \gamma_i)^{\eta_i}
\]
with $\gamma^A > 0$, $\gamma^M = 0$, $\gamma^S < 0$ (so $c^S + \gamma^S$ means $c^S - |\gamma^S|$, i.e., services is a luxury).

\subsection*{Demand Derivation}

The household maximizes $u_t$ subject to the budget constraint $\sum_i p_i c_i = I$.

\begin{derivationbox}[title={Step-by-Step Demand Derivation}]

\textbf{Step 1.} Form the Lagrangian:
\[
\mathcal{L} = \prod_{i} (c_i - \gamma_i)^{\eta_i} - \lambda \!\left(\sum_i p_i c_i - I\right)
\]

\textbf{Step 2.} Take the FOC for good $j$:
\[
\frac{\partial \mathcal{L}}{\partial c_j} = \eta_j (c_j - \gamma_j)^{-1} \cdot \prod_i (c_i - \gamma_i)^{\eta_i} = \lambda p_j
\]
which simplifies to:
\[
\frac{\eta_j\, u}{c_j - \gamma_j} = \lambda p_j \tag{FOC-$j$}
\]

\textbf{Step 3.} Take the ratio of FOC-$j$ and FOC-$i$:
\[
\frac{\eta_j}{\eta_i} \cdot \frac{c_i - \gamma_i}{c_j - \gamma_j} = \frac{p_j}{p_i}
\implies (c_i - \gamma_i) = \frac{\eta_i\, p_j}{\eta_j\, p_i}(c_j - \gamma_j)
\]

\textbf{Step 4.} Substitute into the budget constraint $\sum_i p_i c_i = I$:
\[
\sum_i p_i \gamma_i + \sum_i p_i (c_i - \gamma_i) = I
\]
Using Step 3, $p_i(c_i - \gamma_i) = \eta_i \cdot (c_j - \gamma_j) \cdot p_j / \eta_j$ for all $i$. Summing:
\[
\sum_i p_i (c_i - \gamma_i) = \frac{p_j(c_j - \gamma_j)}{\eta_j} \sum_i \eta_i = \frac{p_j(c_j - \gamma_j)}{\eta_j}
\]
Therefore:
\[
\sum_i p_i \gamma_i + \frac{p_j(c_j - \gamma_j)}{\eta_j} = I
\implies c_j - \gamma_j = \frac{\eta_j (I - \sum_i p_i \gamma_i)}{p_j}
\]

\textbf{Step 5.} Final demand schedule:
\[
\boxed{c_j = \gamma_j + \frac{\eta_j \bigl(I - \sum_i p_i \gamma_i\bigr)}{p_j}}
\tag{Demand}
\]

\end{derivationbox}

\subsection*{Properties of Stone-Geary Demand}

\begin{resultbox}
\begin{itemize}[nosep]
    \item \textbf{Linear Engel curves} but \emph{not} through the origin (non-homothetic).
    \item \textbf{Agriculture ($\gamma^A > 0$):} $c^A - \gamma^A$ grows proportionally with supernumerary income. The Engel curve is concave $\Rightarrow$ necessity. Income elasticity $\varepsilon_I^A = \eta_A I / (p_A c^A) < 1$ since $p_A c^A > p_A \gamma^A + \eta_A(I - \sum p_i \gamma_i) \cdot (p_A/p_A) = \eta_A I + (1-\eta_A) p_A \gamma^A$ ... simplifying: as $I \to \infty$, $\varepsilon_I^A \to \eta_A < 1$.
    \item \textbf{Manufacturing ($\gamma^M = 0$):} Engel curve passes through origin. $c^M$ grows proportionally $\Rightarrow$ income elasticity $\to \eta_M$ (unit elastic asymptotically).
    \item \textbf{Services ($\gamma^S < 0$, i.e., $-|\gamma^S|$):} $c^S + |\gamma^S|$ grows proportionally. The Engel curve is convex $\Rightarrow$ luxury. As $I \to \infty$, $\varepsilon_I^S \to \eta_S < 1$ but for moderate $I$ the elasticity exceeds 1.
    \item \textbf{Budget share of agriculture:} $p_A c^A / I \to \eta_A$ as $I \to \infty$ and $> \eta_A$ for finite $I$.
    \item \textbf{Budget share of services:} $p_S c^S / I \to \eta_S$ as $I \to \infty$ and $< \eta_S$ for finite $I$.
\end{itemize}
\end{resultbox}

\clearpage

%==============================================================================
\section{Demand Model: CGP and Sectoral Growth Rates}
%==============================================================================

\subsection*{The Constant Growth Path (CGP)}

On the \textbf{Constant Growth Path (CGP)}, aggregate consumption $C_t$ grows at a constant rate $g$, but each sector's consumption grows at a \emph{different} rate.

\begin{keyformulabox}[title={CGP Definition}]
A CGP is a balanced-growth-like path on which:
\begin{enumerate}[nosep]
    \item Aggregate consumption grows at constant rate: $\dot{C}/C = g$.
    \item Sectoral consumptions grow at (generally different) constant rates.
    \item Relative prices remain constant.
\end{enumerate}
\textbf{Existence and uniqueness condition:}
\[
\boxed{\frac{\gamma^A}{B_A} = \frac{\gamma^S}{B_S}}
\]
where $B_i$ is a productivity normalization. This condition pins down the unique CGP.
\end{keyformulabox}

\subsection*{Sectoral Consumption Growth Rates on the CGP}

Let aggregate per-capita consumption be $X_t = B_M C_t$ (growing at rate $g$). Then:

\begin{align*}
\frac{\dot{c}_t^A}{c_t^A} &= g \cdot \frac{c_t^A - \gamma^A}{c_t^A} < g
\tag{Agriculture: grows \emph{slower} than $g$} \\[0.5em]
\frac{\dot{c}_t^M}{c_t^M} &= g
\tag{Manufacturing: grows at $g$} \\[0.5em]
\frac{\dot{c}_t^S}{c_t^S} &= g \cdot \frac{c_t^S + \gamma^S}{c_t^S} > g
\tag{Services: grows \emph{faster} than $g$}
\end{align*}

\begin{derivationbox}[title={Intuition for CGP Growth Rate Differences}]
From the demand schedule $c_j = \gamma_j + \eta_j(I - \sum_i p_i \gamma_i)/p_j$, differentiate with respect to time (holding relative prices constant):
\[
\dot{c}_j = \eta_j \dot{I}/p_j = \eta_j g \cdot X_t / p_j
\]
So:
\[
\frac{\dot{c}_j}{c_j} = \frac{\eta_j g X_t / p_j}{c_j} = g \cdot \frac{\eta_j X_t/p_j}{c_j} = g \cdot \frac{c_j - \gamma_j}{c_j}
\]
\begin{itemize}[nosep]
    \item Agriculture ($\gamma^A > 0$): $(c^A - \gamma^A)/c^A < 1$ $\Rightarrow$ grows below $g$.
    \item Manufacturing ($\gamma^M = 0$): $(c^M - 0)/c^M = 1$ $\Rightarrow$ grows at $g$.
    \item Services ($\gamma^S < 0$): $(c^S + |\gamma^S|)/c^S > 1$ $\Rightarrow$ grows above $g$.
\end{itemize}
As $t \to \infty$, all sectoral growth rates converge to $g$ (since $\gamma_i/c_i \to 0$).
\end{derivationbox}

\begin{resultbox}
On the CGP: agriculture grows below aggregate, manufacturing at aggregate, services above aggregate. The \textbf{agriculture share} in total consumption falls; the \textbf{services share} rises; manufacturing share is approximately stable. This matches the demand-side stylized facts without requiring any relative price changes.
\end{resultbox}

\clearpage

%==============================================================================
\section{Employment Dynamics under Demand-Driven Structural Change}
%==============================================================================

\subsection*{Setup}

Each sector $i$ employs $L_{it}$ workers and produces output $Y_{it} = B_i L_{it} F(k^*, 1) \cdot X_t$ (Cobb-Douglas, on balanced growth path with optimal capital per worker $k^*$). Market clearing in sector $i$:
\[
L_{it} \cdot B_i \cdot F(k^*, 1) \cdot X_t = N_t \cdot c_t^i
\]
where $N_t$ is population and $X_t$ is the aggregate productivity index.

\subsection*{Sectoral Employment Growth Rates}

Differentiating the market clearing condition:

\begin{keyformulabox}[title={Employment Growth Rates on the CGP}]
\begin{align*}
\frac{\dot{L}_t^A}{L_t^A} &= n - \frac{g\, \gamma^A L_t}{B_A L_{At} X_t F(k^*,1)} < n
\tag{Agriculture} \\[0.5em]
\frac{\dot{L}_t^M}{L_t^M} &= n
\tag{Manufacturing} \\[0.5em]
\frac{\dot{L}_t^S}{L_t^S} &= n + \frac{g\, \gamma^S L_t}{B_S L_{St} X_t F(k^*,1)} > n
\tag{Services}
\end{align*}
\end{keyformulabox}

\subsection*{Interpretation}

\begin{itemize}
    \item Agricultural employment grows \emph{slower} than population: the agriculture employment share $L^A/L$ falls over time.
    \item Manufacturing employment grows \emph{at} the population rate: the manufacturing share is approximately stable (in the long run converges to $\eta_M$).
    \item Services employment grows \emph{faster} than population: the services employment share $L^S/L$ rises over time.
    \item In the long run, as $c^i \gg \gamma^i$, all employment shares stabilize at $\eta_i$ values.
\end{itemize}

\subsection*{Prediction for the Agriculture Employment Share}

The agriculture employment share $l^A_t = L^A_t/L_t$ satisfies:
\[
\dot{l}^A_t / l^A_t = \frac{\dot{L}^A_t}{L^A_t} - n = -\frac{g\, \gamma^A L_t}{B_A L_{At} X_t F(k^*,1)} < 0
\]

So $l^A_t$ is monotonically decreasing. This matches the key stylized fact.

\clearpage

%==============================================================================
\section{Supply-Side: Baumol's Cost Disease Framework}
%==============================================================================

\subsection*{Two-Sector CES Economy}

Consider a two-sector model (can be thought of as manufacturing vs services, or labor-intensive vs capital-intensive). The final good is a CES aggregate:

\begin{align*}
Y_t = \left[\omega\, Y_{1t}^{\frac{\varepsilon-1}{\varepsilon}} + (1-\omega)\, Y_{2t}^{\frac{\varepsilon-1}{\varepsilon}}\right]^{\frac{\varepsilon}{\varepsilon-1}}
\tag{CES Aggregator}
\end{align*}

where $\varepsilon > 0$ is the elasticity of substitution, $\omega \in (0,1)$ is the share parameter.

\subsection*{Sectoral Production}

\begin{align*}
Y_{it} = A_{it}\, K_{it}^{\alpha_i}\, L_{it}^{1-\alpha_i}, \qquad i = 1, 2
\tag{Sector Production}
\end{align*}

with sector-specific TFP growth:
\[
\frac{\dot{A}_{it}}{A_{it}} = a_i > 0
\]

\textbf{Key asymmetry assumption:} $\alpha_1 < \alpha_2$

\begin{tcolorbox}[colback=orange!10,colframe=orange3,title=Sector Asymmetry]
Sector 1 is \textbf{more labor-intensive} ($\alpha_1 < \alpha_2$); Sector 2 is \textbf{more capital-intensive}.
This capital-intensity difference is what generates factor reallocation dynamics.
\end{tcolorbox}

\subsection*{Cases for Elasticity of Substitution}

\begin{center}
\begin{tabular}{lcl}
\toprule
\textbf{Case} & \textbf{Parameter Value} & \textbf{Relationship between goods} \\
\midrule
Complements & $\varepsilon < 1$ & Hard to substitute; ``need both'' \\
Unit elastic (Cobb-Douglas) & $\varepsilon = 1$ & Cobb-Douglas aggregator \\
Substitutes & $\varepsilon > 1$ & Easy to substitute one for the other \\
Leontief & $\varepsilon \to 0$ & Perfect complements; fixed proportions \\
\bottomrule
\end{tabular}
\end{center}

\clearpage

%==============================================================================
\section{Supply Model: Capital and Labor Allocation}
%==============================================================================

\subsection*{Capital Allocation}

The fraction of total capital allocated to sector 1:
\begin{align*}
k_t := \frac{K_{1t}}{K_t}
\tag{Capital Share}
\end{align*}

From the profit-maximizing firm's FOC equating marginal products of capital across sectors:

\begin{keyformulabox}[title={Capital Allocation Fraction}]
\[
k_t = \left[1 + \frac{(1-\omega)}{\omega} \cdot \frac{\alpha_2}{\alpha_1} \cdot \left(\frac{Y_{2t}}{Y_{1t}}\right)^{\!\frac{\varepsilon-1}{\varepsilon}}\right]^{-1}
\]
\end{keyformulabox}

\subsection*{Labor Allocation}

The fraction of total labor in sector 1:
\begin{align*}
\lambda_t := \frac{L_{1t}}{L_t}
\tag{Labor Share}
\end{align*}

From equating marginal products of labor:

\begin{keyformulabox}[title={Labor Allocation Fraction}]
\[
\lambda_t = \left[1 + \frac{(1-\alpha_2)}{(1-\alpha_1)} \cdot \frac{\alpha_1}{\alpha_2} \cdot \frac{1 - k_t}{k_t}\right]^{-1}
\]
\end{keyformulabox}

\textbf{Key relationship:} $\lambda_t$ is \textbf{increasing in $k_t$}. Labor and capital move in the \emph{same direction}. If capital flows to sector 1, labor also flows to sector 1.

\begin{derivationbox}[title={Why Labor and Capital Move Together}]
With Cobb-Douglas production in each sector and competitive factor markets, the ratio of labor to capital in each sector is determined by the wage-rental ratio $w/r$, which is common across sectors. A higher $k_t$ (more capital in sector 1) means sector 1 has absorbed more productive inputs, raising its demand for labor too. The elasticity relationship $\partial \lambda_t / \partial k_t > 0$ follows from the cross-derivative of the profit function.
\end{derivationbox}

\subsection*{TFP-Driven Reallocation}

The direction of TFP-driven factor flows:
\begin{align*}
\frac{d \ln k_t}{d \ln A_{2t}} &= -\frac{d \ln k_t}{d \ln A_{1t}} > 0 \quad \Longleftrightarrow \quad \varepsilon < 1
\tag{TFP Channel}
\end{align*}

Interpretation: when $\varepsilon < 1$, a TFP increase in sector 2 \emph{draws capital away from sector 2 and toward sector 1}. Counterintuitive, but rational: the fast-growing sector needs fewer inputs to produce its required output.

\clearpage

%==============================================================================
\section{Key Result: Reallocation Direction ($\varepsilon < 1$ vs $\varepsilon > 1$)}
%==============================================================================

\subsection*{Main Theorem}

\begin{keyformulabox}[title={Sign of Capital Reallocation}]
\[
\frac{d \ln k_t}{d \ln K_t} > 0 \quad \Longleftrightarrow \quad (\alpha_2 - \alpha_1)(1 - \varepsilon) > 0
\]
Since $\alpha_2 > \alpha_1$ by assumption, this simplifies to:
\[
\frac{d \ln k_t}{d \ln K_t} > 0 \quad \Longleftrightarrow \quad \varepsilon < 1 \quad \text{(goods are complements)}
\]
\end{keyformulabox}

\subsection*{Two Cases}

\begin{center}
\begin{tabular}{p{3.5cm}p{5.5cm}p{5.5cm}}
\toprule
\textbf{} & \textbf{Complements ($\varepsilon < 1$)} & \textbf{Substitutes ($\varepsilon > 1$)} \\
\midrule
Capital reallocation & Capital flows toward sector 1 (labor-intensive) as $K$ grows & Capital flows toward sector 2 (capital-intensive) as $K$ grows \\
Matching stylized fact & Matches: labor shifts from labor-intensive agriculture to capital-intensive manufacturing & Does not match observed reallocation \\
Intuition & Need both goods; producing more of the scarce good requires the labor-intensive sector to grow & Substitute away from expensive-factor good; capital stays where it has comparative advantage \\
Factor flows & Capital and labor move together toward sector 1 & Capital and labor stay in sector 2 \\
\bottomrule
\end{tabular}
\end{center}

\subsection*{Baumol's Cost Disease: Key Equation}

From the growth-rate accounting of the capital allocation:

\begin{keyformulabox}[title={Baumol Reallocation Condition}]
\begin{align*}
\gamma_{K1} - \gamma_{K2} &= \frac{\varepsilon - 1}{\varepsilon}\, (\gamma_{Y1} - \gamma_{Y2})
\tag{Baumol}
\end{align*}
\end{keyformulabox}

Since labor follows capital ($\lambda_t$ is increasing in $k_t$), the same sign applies to labor:

\begin{resultbox}
\textbf{If $\varepsilon < 1$ (complements):}
\[
\gamma_{L1} > \gamma_{L2} \iff \gamma_{K1} > \gamma_{K2} \iff \gamma_{Y1} < \gamma_{Y2}
\]
The \textbf{slow-growing sector absorbs factors}. Reallocation goes \emph{opposite} to sectoral productivity growth.\\[0.4em]
\textbf{If $\varepsilon > 1$ (substitutes):}
\[
\gamma_{L1} < \gamma_{L2} \iff \gamma_{K1} < \gamma_{K2} \iff \gamma_{Y1} < \gamma_{Y2}
\]
The \textbf{fast-growing sector absorbs factors}. Reallocation goes \emph{in the same direction} as productivity growth.
\end{resultbox}

\subsection*{Baumol's Cost Disease Interpretation}

Services have \emph{low} TFP growth. If services and manufacturing are complements ($\varepsilon < 1$):
\begin{itemize}
    \item Services become relatively expensive (price rises relative to manufacturing).
    \item But demand is inelastic, so nominal spending on services rises.
    \item Factors (labor) flow \emph{into} services, even though services are becoming more expensive.
    \item This is the ``cost disease'': rising costs in services do not reduce demand due to complementarity.
\end{itemize}

\begin{tcolorbox}[colback=orange!10,colframe=orange3,title=Key Insight to Memorize]
\textbf{The Baumol insight:} When $\varepsilon < 1$, the low-TFP-growth sector (services) absorbs factors over time, even as it becomes relatively more expensive. This explains the services employment share rising monotonically.
\end{tcolorbox}

\clearpage

%==============================================================================
\section{Testing Demand vs Supply: Herrendorf, Rogerson \& Valentinyi (2013)}
%==============================================================================

\subsection*{Why the Test Matters}

Both demand and supply mechanisms can in principle generate the observed structural change patterns. Herrendorf, Rogerson \& Valentinyi (HRV) ask: which mechanism quantitatively dominates, and can we identify the structural parameters?

\subsection*{HRV Preferences}

HRV use the following non-homothetic CES preferences:
\begin{align*}
u_t = \left[\sum_{i \in \{a,m,s\}} \omega_i^{1/\sigma}\,(c_{it} - \bar{c}_i)^{\frac{\sigma-1}{\sigma}}\right]^{\frac{\sigma}{\sigma-1}}
\tag{HRV Utility}
\end{align*}

where:
\begin{itemize}[nosep]
    \item $\sigma > 0$: intertemporal/intratemporal elasticity of substitution.
    \item $\omega_i$: weight on sector $i$ in the aggregator.
    \item $\bar{c}_i$: non-homothetic threshold (positive $\Rightarrow$ necessity; negative $\Rightarrow$ luxury).
\end{itemize}

\subsection*{Two Approaches to Measurement}

\begin{center}
\begin{tabular}{p{4cm}p{5cm}p{5cm}}
\toprule
\textbf{} & \textbf{Final Expenditure Approach} & \textbf{Value Added Approach} \\
\midrule
What is measured & What households actually buy (e.g., food in restaurants, movies, cars) & Value added in each sector's production (agriculture, manufacturing, services as classified by GDP accounts) \\
What substitution is & Substituting between consumer goods (e.g., eating out vs grocery shopping) & Substituting between bundles of factors organized by industry \\
\bottomrule
\end{tabular}
\end{center}

\subsection*{HRV Empirical Results}

\begin{center}
\begin{tabular}{lcc}
\toprule
\textbf{Parameter} & \textbf{Final Expenditure} & \textbf{Value Added} \\
\midrule
$\sigma$ (elasticity of substitution) & 0.85 & $\approx 0.002$ (nearly Leontief!) \\
$\bar{c}_a$ (agriculture threshold) & $-1350.38$ (luxury!) & $-138.68$ \\
$\bar{c}_s$ (services threshold) & $11{,}237.40$ (necessity) & $4{,}261.82$ \\
$\omega_s$ (services weight) & $0.81$ & --- \\
\bottomrule
\end{tabular}
\end{center}

\subsection*{Interpretation of the Divergence}

\begin{derivationbox}[title={Why Do the Two Approaches Disagree?}]
\textbf{Final expenditure:} When a consumer chooses between watching a movie (services) and attending a football match (also services), or between a restaurant meal and a supermarket meal, these are \emph{genuine substitutes} from the consumer's perspective. Hence $\sigma \approx 0.85$: moderate substitutability.

\textbf{Value added:} A restaurant meal and supermarket food use \emph{very similar factor inputs} (labor, capital). In GDP accounts, a restaurant meal is classified as \emph{services} but supermarket food as \emph{agriculture}. Since their factor compositions are almost identical, the value-added aggregator has essentially Leontief complementarity: $\sigma \approx 0.002$.

\textbf{Implication:} The ``structural change'' seen in value-added sector shares partly reflects how national accounts classify goods, not genuine consumer substitution. This matters for identifying the demand vs supply mechanism.
\end{derivationbox}

\subsection*{Classification Puzzle: Agriculture as a Luxury}

A striking HRV result is that agriculture has $\bar{c}_a < 0$ in the final expenditure approach, making it a \textbf{luxury}. This seems to contradict Engel's Law. Why?

\begin{resultbox}
In the \textbf{final expenditure} classification, what counts as ``agriculture'' includes restaurant meals, prepared foods, premium ingredients — goods that are luxuries. Low-income consumers eat cheap staples (classified as agriculture); high-income consumers eat processed foods and restaurant meals (classified as services). So in final expenditure terms, \emph{agriculture spending} can rise more than proportionally with income.

In the \textbf{value added} classification, agriculture's $\bar{c}_a < 0$ still holds but is less extreme, because the distinction is by production process (farming vs manufacturing vs services).
\end{resultbox}

\clearpage

%==============================================================================
\section{Empirical Results and Literature}
%==============================================================================

\subsection*{Summary of Key Papers}

\begin{center}
\begin{tabular}{p{3.5cm}p{5cm}p{5.5cm}}
\toprule
\textbf{Paper} & \textbf{Mechanism / Contribution} & \textbf{Key Finding / Limitation} \\
\midrule
Duarte \& Restuccia (2010) & Documents TFP growth differentials across sectors 1956--2004 & Agriculture 3.8\%, Manufacturing 2.4\%, Services 1.3\%; TFP gaps explain much of cross-country income differences \\[0.5em]
Herrendorf, Rogerson \& Valentinyi (2013) & Tests demand vs supply using two measurement approaches & Final expenditure: $\sigma \approx 0.85$, demand matters. Value added: $\sigma \approx 0$, supply (Baumol) dominates; reconciles the two views \\[0.5em]
Buera \& Kaboski (2009) & Stone-Geary + sector-biased technical change & Cannot match the steep manufacturing employment decline at high income levels \\[0.5em]
Boppart (2014) & PIGL (Price-Independent Generalized Linear) preferences & Non-homotheticities do not vanish asymptotically; demand and supply each explain roughly 50\% of structural change \\[0.5em]
Comin, Lashkari \& Mestieri (2021) & Non-homothetic CES preferences & $R^2$ twice as large as simpler models; non-homotheticities explain approximately 3/4 of structural change \\
\bottomrule
\end{tabular}
\end{center}

\subsection*{Demand vs Supply: Side-by-Side Comparison}

\begin{center}
\begin{tabular}{p{3.5cm}p{5.5cm}p{5.5cm}}
\toprule
\textbf{Feature} & \textbf{Demand (Non-Homotheticity)} & \textbf{Supply (Baumol Cost Disease)} \\
\midrule
Key mechanism & Income effects: as income rises, budget shares shift toward luxuries (services) and away from necessities (agriculture) & Price effects: productivity differences change relative prices; with $\varepsilon < 1$, factors flow toward slow-TFP sector \\
Key parameter & Income elasticities ($\bar{c}_i$ in Stone-Geary / HRV) & Elasticity of substitution $\varepsilon$ \\
Relative prices & Structural change occurs even with \emph{constant} relative prices & Relative price changes drive factor reallocation \\
Cross-country variation & Predicts SC correlated with per-capita income level & Predicts SC driven by TFP growth differentials \\
Long-run behavior & Non-homotheticities vanish as $I \to \infty$ (sector shares converge to $\eta_i$) & Persists as long as TFP growth rates differ \\
Evidence & HRV: $\sigma \approx 0.85$ (final expenditure); non-homothetic CES fits well & HRV: $\sigma \approx 0$ (value added); Baumol mechanism important \\
Limitation & Cannot explain all of manufacturing's inverse-U pattern & With $\varepsilon > 1$, predicts wrong direction of reallocation \\
Best suited for & Cross-section: poor vs rich countries & Time-series: explaining secular increase in services share \\
\bottomrule
\end{tabular}
\end{center}

\clearpage

%==============================================================================
\section{Problem-Solving Templates}
%==============================================================================

\subsection*{(a) Derive Stone-Geary Demand Schedules}

\begin{enumerate}
    \item Write the Stone-Geary utility: $u = \prod_i (c_i - \gamma_i)^{\eta_i}$.
    \item Set up the Lagrangian with budget constraint $\sum_i p_i c_i = I$.
    \item Derive FOC for good $j$: $\eta_j u / (c_j - \gamma_j) = \lambda p_j$.
    \item Take ratios of FOCs to eliminate $\lambda$: $\eta_i (c_j - \gamma_j) p_j = \eta_j (c_i - \gamma_i) p_i$.
    \item Substitute into budget constraint. Use $\sum_i \eta_i = 1$ to simplify.
    \item Result: $c_j = \gamma_j + \eta_j (I - \sum_i p_i \gamma_i) / p_j$.
    \item Verify: check that $\sum_i p_i c_i = I$ holds as a consistency check.
\end{enumerate}

\subsection*{(b) Compute Engel Curve Slope and Income Elasticity}

\begin{enumerate}
    \item From demand $c_j = \gamma_j + \eta_j(I - \sum_k p_k \gamma_k)/p_j$, differentiate w.r.t. $I$:
    \[
    \frac{\partial c_j}{\partial I} = \frac{\eta_j}{p_j} > 0 \quad \text{(always a normal good)}
    \]
    \item Income elasticity:
    \[
    \varepsilon_I^j = \frac{\partial c_j}{\partial I} \cdot \frac{I}{c_j} = \frac{\eta_j I}{p_j c_j}
    \]
    \item Substitute $p_j c_j = p_j \gamma_j + \eta_j (I - \sum_k p_k \gamma_k)$:
    \[
    \varepsilon_I^j = \frac{\eta_j I}{p_j \gamma_j + \eta_j (I - \sum_k p_k \gamma_k)}
    \]
    \item Sign of $\varepsilon_I^j - 1$: positive (luxury) iff $p_j \gamma_j < 0$ (i.e., $\gamma_j < 0$); negative (necessity) iff $\gamma_j > 0$.
    \item As $I \to \infty$: $\varepsilon_I^j \to \eta_j < 1$ for agriculture (since $\gamma^A > 0$ is finite and $I$ dominates).
\end{enumerate}

\subsection*{(c) Derive CGP Growth Rates}

\begin{enumerate}
    \item Start from demand: $c_j = \gamma_j + \eta_j (I - \sum_k p_k \gamma_k)/p_j$.
    \item Differentiate with respect to time (holding prices constant): $\dot{c}_j = \eta_j \dot{I}/p_j$.
    \item Identify $\dot{I}/p_j$ in terms of $c_j - \gamma_j$: from demand, $(c_j - \gamma_j) = \eta_j(I - \sum p_k \gamma_k)/p_j$, so $\eta_j \dot{I}/p_j = \dot{I}/(I - \sum p_k \gamma_k) \cdot (c_j - \gamma_j)$.
    \item Define $g = \dot{I}/(I - \sum p_k \gamma_k)$ as the aggregate growth rate of supernumerary income.
    \item Sector growth rate: $\dot{c}_j/c_j = g (c_j - \gamma_j)/c_j$.
    \item For agriculture ($\gamma^A > 0$): ratio $< 1$ $\Rightarrow$ grows below $g$.
    \item For services ($\gamma^S < 0$): ratio $> 1$ $\Rightarrow$ grows above $g$.
\end{enumerate}

\subsection*{(d) Apply the Baumol Reallocation Condition}

\begin{enumerate}
    \item Identify which sector is more capital-intensive: $\alpha_2 > \alpha_1$ (sector 2 capital-intensive).
    \item Identify elasticity of substitution $\varepsilon$ from the problem.
    \item Apply the sign rule:
    \begin{itemize}
        \item $\varepsilon < 1$: $d\ln k_t / d\ln K_t > 0$ $\Rightarrow$ capital flows toward sector 1 (labor-intensive).
        \item $\varepsilon > 1$: $d\ln k_t / d\ln K_t < 0$ $\Rightarrow$ capital flows toward sector 2 (capital-intensive).
    \end{itemize}
    \item Use the Baumol equation: $\gamma_{K1} - \gamma_{K2} = [(\varepsilon-1)/\varepsilon](\gamma_{Y1} - \gamma_{Y2})$.
    \item Since $\lambda_t$ is increasing in $k_t$: labor and capital flow in the same direction.
    \item Conclude on the direction of labor reallocation and which sector's employment share rises.
\end{enumerate}

\subsection*{(e) Identify Demand vs Supply Channel (Empirical Strategy)}

\begin{enumerate}
    \item Compare value-added vs final expenditure measures of sector shares and relative prices.
    \item Estimate $\sigma$ / $\varepsilon$ from each approach separately.
    \item \textbf{If $\sigma$ (final expenditure) is substantially higher than $\sigma$ (value added)}: evidence that classification matters and supply mechanism is important in value added.
    \item Check $\bar{c}_i$ signs: negative $\bar{c}_a$ in final expenditure $\Rightarrow$ agriculture classified as luxury in consumer spending (HRV finding).
    \item Use $\sigma \approx 0$ (value added) to argue Baumol/supply dominates in national accounts.
    \item Use $\sigma \approx 0.85$ (final expenditure) to argue demand substitution matters for consumer allocation.
\end{enumerate}

\clearpage

%==============================================================================
\section{Common Pitfalls}
%==============================================================================

\begin{enumerate}

    \item \textbf{Confusing homotheticity with unit income elasticity everywhere}
    \begin{itemize}
        \item Homothetic preferences have income elasticity of 1 for all goods, so all Engel curves pass through the origin.
        \item Stone-Geary is \emph{not} homothetic: income elasticities differ across goods and depend on $I$.
        \item Do not say ``Stone-Geary is homothetic because it has linear Engel curves'': linearity without passing through origin is not homotheticity.
    \end{itemize}

    \item \textbf{Forgetting the sign of $\gamma^S$ in Stone-Geary}
    \begin{itemize}
        \item The standard formulation is $u = (c^A - \gamma^A)^{\eta_A}(c^M)^{\eta_M}(c^S + \gamma^S)^{\eta_S}$ with $\gamma^A, \gamma^S > 0$.
        \item The $+\gamma^S$ means services acts as a luxury (negative threshold in the unified notation $c_i - \gamma_i$ with $\gamma^S_{\text{unified}} < 0$).
        \item Always verify: $c^A - \gamma^A > 0$ must hold (consumption above subsistence).
    \end{itemize}

    \item \textbf{Misidentifying the direction of Baumol reallocation}
    \begin{itemize}
        \item The key result is: with $\varepsilon < 1$, factors flow toward the \emph{slow-TFP-growth} sector.
        \item Many students incorrectly say ``factors follow productivity''. The opposite is true under complementarity.
        \item Mnemonic: ``slow sector starves for inputs because you cannot substitute away from it''.
    \end{itemize}

    \item \textbf{Applying $\varepsilon > 1$ intuition when the model uses $\varepsilon < 1$}
    \begin{itemize}
        \item For structural change to match stylized facts in the supply model, we need $\varepsilon < 1$.
        \item Do not apply standard comparative advantage arguments (factors flow to productive sector) --- that reasoning applies when $\varepsilon > 1$ only.
    \end{itemize}

    \item \textbf{Confusing the CGP with a balanced growth path}
    \begin{itemize}
        \item On a BGP, all variables grow at the same rate. On the CGP, sector growth rates differ.
        \item The CGP is the correct concept for structural change: it allows sector heterogeneity while keeping aggregate growth constant.
        \item Existence of the CGP requires the specific condition $\gamma^A / B_A = \gamma^S / B_S$.
    \end{itemize}

    \item \textbf{Conflating final expenditure and value added in empirical analysis}
    \begin{itemize}
        \item HRV's key insight is that final expenditure and value added give very different $\sigma$ estimates.
        \item Do not assume they should agree; their divergence reveals information about the classification of goods.
        \item When told ``services and manufacturing are complements'', check whether this refers to final goods or factor-input bundles.
    \end{itemize}

    \item \textbf{Claiming structural change requires relative price changes}
    \begin{itemize}
        \item With demand non-homotheticities (Stone-Geary), structural change occurs even at \emph{constant relative prices}.
        \item Income growth alone shifts budget shares if income elasticities differ.
        \item Relative price changes are necessary only for the supply (Baumol) mechanism.
    \end{itemize}

    \item \textbf{Ignoring that labor and capital move together}
    \begin{itemize}
        \item In the supply model, $\lambda_t$ is increasing in $k_t$, so labor always follows capital.
        \item Once you determine which sector receives more capital, the labor allocation is determined.
    \end{itemize}

\end{enumerate}

\clearpage

%==============================================================================
\section{Quick Reference Formulas}
%==============================================================================

\subsection*{Stone-Geary Preferences}
\begin{align*}
u &= (c^A - \gamma^A)^{\eta_A} (c^M)^{\eta_M} (c^S + \gamma^S)^{\eta_S}, \quad \eta_A + \eta_M + \eta_S = 1 \tag{Utility}\\
c_j &= \gamma_j + \frac{\eta_j (I - \sum_i p_i \gamma_i)}{p_j} \tag{Demand}\\
\frac{\partial c_j}{\partial I} &= \frac{\eta_j}{p_j}, \qquad \varepsilon_I^j = \frac{\eta_j I}{p_j c_j} \tag{Engel Curve}
\end{align*}

\subsection*{CGP Growth Rates}
\begin{align*}
\frac{\dot{c}^A}{c^A} &= g \cdot \frac{c^A - \gamma^A}{c^A} < g \tag{Agriculture}\\
\frac{\dot{c}^M}{c^M} &= g \tag{Manufacturing}\\
\frac{\dot{c}^S}{c^S} &= g \cdot \frac{c^S + \gamma^S}{c^S} > g \tag{Services}
\end{align*}

\subsection*{Employment Growth Rates (Demand Model)}
\begin{align*}
\frac{\dot{L}^A}{L^A} &= n - \frac{g\, \gamma^A L}{B_A L^A X F(k^*,1)} < n \tag{Agriculture}\\
\frac{\dot{L}^M}{L^M} &= n \tag{Manufacturing}\\
\frac{\dot{L}^S}{L^S} &= n + \frac{g\, \gamma^S L}{B_S L^S X F(k^*,1)} > n \tag{Services}
\end{align*}

\subsection*{Supply Model: CES Aggregator and Allocation}
\begin{align*}
Y_t &= \left[\omega Y_{1t}^{\frac{\varepsilon-1}{\varepsilon}} + (1-\omega) Y_{2t}^{\frac{\varepsilon-1}{\varepsilon}}\right]^{\frac{\varepsilon}{\varepsilon-1}} \tag{CES}\\
k_t &= \left[1 + \frac{1-\omega}{\omega} \cdot \frac{\alpha_2}{\alpha_1} \cdot \left(\frac{Y_{2t}}{Y_{1t}}\right)^{\!\frac{\varepsilon-1}{\varepsilon}}\right]^{-1} \tag{Capital Allocation}\\
\lambda_t &= \left[1 + \frac{1-\alpha_2}{1-\alpha_1} \cdot \frac{\alpha_1}{\alpha_2} \cdot \frac{1-k_t}{k_t}\right]^{-1}, \quad \frac{d\lambda_t}{dk_t} > 0 \tag{Labor Allocation}
\end{align*}

\subsection*{Key Reallocation Results}
\begin{align*}
\frac{d\ln k_t}{d\ln K_t} &> 0 \iff (\alpha_2 - \alpha_1)(1 - \varepsilon) > 0 \iff \varepsilon < 1 \tag{Capital Direction}\\
\gamma_{K1} - \gamma_{K2} &= \frac{\varepsilon-1}{\varepsilon}\,(\gamma_{Y1} - \gamma_{Y2}) \tag{Baumol}\\
\frac{d\ln k_t}{d\ln A_{2t}} &> 0 \iff \varepsilon < 1 \tag{TFP Channel}
\end{align*}

\subsection*{HRV Utility and Empirical Estimates}
\begin{align*}
u_t &= \left[\sum_{i} \omega_i^{1/\sigma}(c_{it} - \bar{c}_i)^{\frac{\sigma-1}{\sigma}}\right]^{\frac{\sigma}{\sigma-1}} \tag{HRV}
\end{align*}

\begin{center}
\begin{tabular}{lccc}
\toprule
\textbf{Parameter} & \textbf{Final Expenditure} & \textbf{Value Added} & \textbf{Implication} \\
\midrule
$\sigma$ & 0.85 & 0.002 & Leontief in VA; substitutable in FE \\
$\bar{c}_a$ & $-1350.38$ & $-138.68$ & Agriculture is luxury in FE \\
$\bar{c}_s$ & $11{,}237.40$ & $4{,}261.82$ & Services is necessity \\
\bottomrule
\end{tabular}
\end{center}

\subsection*{Classification of Goods (Quick Reference)}

\begin{center}
\begin{tabular}{lcc}
\toprule
\textbf{Sector} & \textbf{Income Elasticity} & \textbf{$\gamma_i$ Sign (Stone-Geary)} \\
\midrule
Agriculture & $<1$ (necessity; $\varepsilon_I^A \to \eta_A$ as $I \to \infty$) & $\gamma^A > 0$ \\
Manufacturing & $\approx 1$ & $\gamma^M = 0$ \\
Services & $>1$ for moderate $I$ (luxury) & $\gamma^S < 0$ (i.e., $+\gamma^S$ in utility with $\gamma^S > 0$) \\
\bottomrule
\end{tabular}
\end{center}

\end{document}
